\subsection*{一步一步建立直观}

\paragraph{1. 先辨认信号是在衰减还是增长}

右图的二维轨线也可以看成复信号 $z(t)=e^{(\lambda+i\omega)t}z(0)$。选择衰减情形，幅度逐渐减小。拉普拉斯变换给信号再乘一个指数权重后积分，让我们用一个复参数同时探测衰减与振荡，而不仅是读取某个时刻的值。



\paragraph{2. 探测频率与信号频率相互比较}

纯旋转信号幅度不衰减；若只无限累计它，普通意义的收敛就需要格外小心。乘上探测核 $e^{-st}$ 后，$s$ 的虚部比较旋转频率，实部提供衰减。右图展示被探测的信号，不直接画积分；应把信号轨迹与变换结果这两个对象区分开。



\paragraph{3. 收敛域是计算的一部分}

改成增长信号。若其增长率为 $\lambda$，要使指数加权后的绝对积分收敛，需要 $\operatorname{Re}s>\lambda$。这一条件说明“给函数写一个变换式”不能脱离参数范围。图像可以提示增长趋势，最终的收敛判断仍由积分估计完成。
